\begin{table*}[t]
  \centering
  \caption{五段固定的 150 帧 Harmony4D 多镜头视频流上的完整系统运行效率（共 750 帧）。}
  \label{tab:end-to-end-runtime-zh}
  \small
  \setlength{\tabcolsep}{8pt}
  \renewcommand{\arraystretch}{0.96}
  \begin{tabular}{llccrr}
    \toprule
    类别 & 方法 & 相机 & 场景 &
    秒 / 150 帧 $\downarrow$ & FPS $\uparrow$ \\
    \midrule
    \multirow{3}{*}{离线}
      & PromptHMR (SPEC) & \xmark & \xmark & $708.7 \pm 209.0$ & 0.212 \\
      & TRAM & \cmark & \xmark & $455.5 \pm 60.6$ & 0.329 \\
      & JOSH & \cmark & \cmark & $949.6 \pm 286.0$ & 0.158 \\
    \midrule
    半在线
      & OnlineHMR & \cmark & \xmark & $1007.8 \pm 124.7$ & 0.149 \\
    \midrule
    \multirow{3}{*}{在线}
      & TRACE & \xmark & \xmark & $45.0 \pm 21.5$ & 3.334 \\
      & \humanthree{} & \cmark & \cmark & $120.4 \pm 3.2$ & 1.246 \\
      & \textbf{\method{}} & \cmark & \cmark & $167.8 \pm 4.0$ & 0.894 \\
    \bottomrule
  \end{tabular}
  \vspace{0.5mm}

  \parbox{0.91\textwidth}{\footnotesize
  秒数为均值 $\pm$ 总体标准差，FPS 为 750 除以总墙钟时间。相机表示流程输出可独立评测的物理相机轨迹，场景表示输出场景几何。计时覆盖方法所需的预测进程启动和模型加载，直至原生预测写盘；多阶段流程报告必要阶段之和，不包含渲染和评测。因此，解释吞吐率时仍需结合输出范围。}
\end{table*}

